% !TEX program = xelatex
% =====================================================================
% Compile with XeLaTeX (needed for the Chinese characters).
% On Overleaf: Menu -> Compiler -> XeLaTeX.
% =====================================================================
\documentclass[10pt,twocolumn]{article}

\usepackage[a4paper,margin=0.85in]{geometry}
\setlength{\columnsep}{0.62cm}
\usepackage{fontspec}
\usepackage{xeCJK}
% Robust CJK font selection: prefer Noto Serif CJK TC, fall back to Fandol.
\IfFontExistsTF{Noto Serif CJK TC}%
  {\setCJKmainfont{Noto Serif CJK TC}}%
  {\IfFontExistsTF{FandolSong}{\setCJKmainfont{FandolSong}}{}}

\usepackage{microtype}
\usepackage{amsmath}
\usepackage{booktabs}
\usepackage{tabularx}
\usepackage{siunitx}
\usepackage{graphicx}
\usepackage{caption}
\captionsetup{font=small,labelfont=bf}
\usepackage[hidelinks]{hyperref}
\usepackage{xcolor}
\hypersetup{colorlinks=true,linkcolor=black,citecolor=black,urlcolor=blue!55!black}
\usepackage{titlesec}
\titleformat{\section}{\large\bfseries}{\thesection}{0.6em}{}
\titleformat{\subsection}{\normalsize\bfseries}{\thesubsection}{0.6em}{}
\titlespacing*{\section}{0pt}{1.4ex plus .3ex}{0.8ex}
\setlength{\parskip}{0.35em}
\sisetup{detect-weight=true}

\newcommand{\dens}[1]{\mbox{#1}}

\title{\bfseries The Half-Life of a Lost Empire:\\[2pt]
\large Measuring the Decay of Loyalist Allusion in Taiwan's\\
Classical Poetry across the 1895 Rupture}
\author{Aryan Maity\\[2pt]
St.~Edmund's College\\[2pt]
\normalsize\href{https://github.com/aryan35790jp/taiwan-loyalist-allusion}{\texttt{github.com/aryan35790jp/taiwan-loyalist-allusion}}}
\date{}

\begin{document}
\makeatletter
\twocolumn[
  \begin{@twocolumnfalse}
  \maketitle

\begin{abstract}
\noindent
When the Qing ceded Taiwan to Japan in 1895, the island's classical poets faced a second
fall of a realm within living cultural memory of the first, the Ming collapse that had once
driven loyalists across the strait to Koxinga's Taiwan. This paper asks whether that second
rupture left a measurable trace in the language of the verse itself. Working only from
openly available texts and an auditable lexicon of dynastic--loss and loyalist allusion
(\mbox{黍離}, \mbox{採薇}, \mbox{文天祥}, \mbox{屈原}, and the local \mbox{延平}--\mbox{東寧}
tradition), we score the density of loyalist allusion across a corpus dated by era and find
a clear bimodal shape: the register is loud at the Ming--Zheng founding, falls almost silent
through the settled mid-Qing, and rises again into the cession years. Against the stable
mid-Qing baseline the rupture-era density is roughly three times higher
($0.087 \to 0.268$ per hundred characters, $p<0.0001$), while a frequency-matched set of
neutral canonical allusions stays flat, so the rise is specific to loyalism rather than a
general taste for allusion. The surge is carried by the classical canon of martyrdom and
refusal-to-serve ($\mbox{採薇}/\mbox{文天祥}/\mbox{田橫}/\mbox{蘇武}$, $p=0.0065$), not by the
local Koxinga vocabulary, and it survives every variation of the detector we tried. Two
independent checks support the reading: an officially compiled gazetteer corpus reproduces
the mid-Qing baseline almost exactly ($p=0.71$), and a classical-Chinese language model that
never saw our lexicon nonetheless separates loyalist from control poems well above chance
($0.70$ vs.\ $0.51$, $p<0.0001$), which argues that the register is a real property of the
language and not an artefact of a word list. We are candid about what the present data cannot
settle: the single corpus that would close the question of editorial selection at the rupture
peak is a national database that is no longer reachable, and we mark that analysis as the
clear next step rather than claim it. What remains is a quantified, falsifiable account of a
literature reaching for an inherited language of loss at the moment its world changed.
\end{abstract}
  \vspace{1.2em}
  \end{@twocolumnfalse}
]
\makeatother

\section{Introduction}
The standard history of Taiwan's classical poetry is usually told as a story of communities:
who founded which poetry society, which anthologies preserved whose work, and how the
literati negotiated life first under the Qing and then under Japanese rule. Running through
that history is a quieter thread that scholars of the period have long recognised but rarely
measured, namely a recurring language of dynastic loss. The poets who followed Koxinga to
Taiwan after the Ming fell wrote as \mbox{遺民}, ``leftover subjects'' of a house that no
longer ruled, and they reached for a fixed repertoire of classical figures, the millet of
the ruined Zhou capital, the brothers who starved rather than eat Zhou grain, the loyal
minister drowned in the Xiang, to say what could not be said directly. Two and a half
centuries later, when the Treaty of Shimonoseki handed the island to Japan in 1895, a new
generation of Taiwanese poets reached, it has often been observed, for the same repertoire.

The observation is old; the measurement is not. Literary historians have read the loyalist
strain in colonial-era Taiwan poetry closely and persuasively, but case by case, poet by
poet. What no one has done, to our knowledge, is treat the repertoire as a countable object
and ask how its frequency moves across the whole span of the tradition. That is the question
of this paper. If the second fall of a realm reactivated the first's language of mourning,
the trace should be visible not only in the poems historians already cite but in the
aggregate: loyalist allusion should be measurably denser in the verse written around 1895
than in the verse of the long, settled century before it, and the rise should be specific to
loyalism rather than a by-product of poets simply alluding more.

We answer this with a deliberately conservative method. From authoritative reference sources
we build a graded lexicon of dynastic-loss and loyalist allusion, every entry anchored to a
classical locus and weighted by how exclusively it signals loyalism. We match it, by exact
and human-checkable string matching, against poems dated by era, and we measure the density
of loyalist allusion per hundred characters against a frequency-matched set of neutral
canonical allusions that serves as a control. The headline result is a bimodal curve: the
loyalist register is loud at the Ming--Zheng founding, nearly silent in the mid-Qing, and
loud again at the cession. Our contributions are:

\begin{itemize}\setlength{\itemsep}{2pt}
\item a graded, openly auditable lexicon of loyalist allusion for classical Chinese, with a
frequency-matched neutral control, designed so that every measurement can be checked by hand;
\item a bimodal diachronic finding---loyalist allusion peaking at both of Taiwan's sovereignty
ruptures and falling quiet between them---that is significant, specific to loyalism, and
robust to the matching choices;
\item a decomposition showing that the 1895 surge is driven by the classical canon of
martyrdom and refusal-to-serve rather than by the local Koxinga vocabulary;
\item two independent corroborations (an officially compiled gazetteer baseline and a
language-model probe that never saw the lexicon); and
\item an honest account of the one question the available data cannot close, and a
pipeline ready to close it the moment the national corpus becomes reachable again.
\end{itemize}

\section{Background}
\subsection{Loyalism and the language of dynastic loss}
The figure of the \mbox{遺民}, the subject who outlives the dynasty he served and refuses to
serve its successor, is one of the oldest in the Chinese literary tradition, and it carries
with it a settled vocabulary. The mourning of a fallen capital is spoken through
\mbox{黍離}, the millet growing over the ruins of Zhou; refusal to take office under the
conqueror through Bo Yi and Shu Qi, who starved on Mount Shouyang rather than eat the grain
of Zhou (\mbox{採薇}); loyal death through Wen Tianxiang, the Song minister executed by the
Yuan, and through the five hundred followers of Tian Heng who died together rather than
submit. Taiwan's place in this tradition is unusually concrete, because the island was itself
a loyalist refuge: Koxinga (\mbox{鄭成功}, Prince of \mbox{延平}) held it for the Ming cause,
and the polity his family ran was named \mbox{東寧}. When 1895 came, the inherited language
of the first fall lay ready to hand for the second, and with a local edge the mainland
tradition did not have.

Scholarship on colonial-era Taiwan literature has read this loyalist strain with care, most
fully in studies of loyalism and the figure of the loyal subject under Japanese rule. That
work is qualitative and interpretive, and it is the foundation we build on; our aim is not to
displace it but to give it a quantitative spine, to ask whether the pattern close reading
finds in individual poets holds across the corpus when counted.

\subsection{Computational study of allusion}
Detecting allusion and intertextual reuse computationally is an established line of work in
the digital humanities, and recent studies have applied it to classical Chinese poetry, both
through phrase matching against a reference canon and through language models trained on
classical text. We take the simplest defensible position available: lexicon matching, in
which every hit is a known figure with a known classical source and can be inspected by a
reader. This trades recall for transparency, which is the right trade when the object of
study is a contested historical claim rather than a benchmark score. To guard against the
obvious worry that a lexicon merely finds itself, we complement it with a representational
probe using a classical-Chinese language model that has no knowledge of our list.

\section{Data}
\subsection{The dated anthology}
Our primary diachronic corpus is \mbox{連橫}'s \mbox{《臺灣詩乘》} (\emph{Taiwan Poetry
Records}, c.\ 1921), a six-volume critical anthology that quotes poems within editorial prose
that narrates the island's literary history in rough chronological order. We extract every
bracketed quoted poem and date it by the nearest preceding temporal marker in the surrounding
prose (a reign era such as \mbox{光緒}, or a dated event such as the cession), which yields
997 poems carrying an approximate year between 1665 and 1896. The procedure uses the
anthology's own chronology and requires no manual dating.

\subsection{Independent and single-author corpora}
Because an anthology reflects its editor's choices, we add two corpora that \mbox{連橫} did
not compile. The first is the complete chronological collection of a single poet who lived
through the rupture, \mbox{許南英}'s \mbox{《窺園留草》} (the poet organised the defence of
Tainan in 1895 and then fled to the mainland); its sections are dated by year, giving an
uncurated within-author series from 1884 to 1917. The second is a set of 867 regulated-verse
poems extracted from official Qing gazetteers (\mbox{諸羅縣志}, \mbox{臺灣縣志},
\mbox{重修臺灣縣志}, \mbox{臺海使槎錄}), compiled by Qing officials rather than by any
loyalist, which provides an independent early/mid-Qing baseline. For a poet-level validation
we also collect the author-attributed poems of \mbox{丘逢甲}, \mbox{連橫}, and the modernist
\mbox{賴和} from a public text repository.

\subsection{Provenance and honesty about access}
All texts above are openly readable. The one corpus that would settle the question of
editorial selection at the rupture peak, the national \mbox{全臺詩} (\emph{Complete Taiwanese
Poems}) database of roughly fifty thousand poems, is no longer reachable online (its host no
longer resolves) and survives in print and behind a commercial subscription. Our pipeline is
written to ingest its export unchanged; until that is possible, we report what the reachable
corpora can and cannot establish, and mark the population analysis as future work rather than
fold an inaccessible result into our claims.

\section{Method}
\subsection{A graded loyalist lexicon and a neutral control}
The loyalist lexicon collects the standard repertoire of dynastic-loss and loyalist allusion
across six families: mourning a fallen realm (\mbox{黍離}, \mbox{銅駝}, \mbox{崖山}),
loyalty and refusal to serve (\mbox{採薇}, \mbox{文天祥}, \mbox{田橫}, \mbox{蘇武}), the
exiled loyal minister (\mbox{屈原}, \mbox{庾信}, \mbox{新亭}), withdrawal as political refusal
(\mbox{桃源}, \mbox{商山四皓}), the local Ming--Zheng layer (\mbox{延平}, \mbox{東寧},
\mbox{寧靖王}, \mbox{五妃}), and explicit loyalist diction (\mbox{遺民}, \mbox{孤臣},
\mbox{割臺}). Each entry lists its surface forms, is anchored to a classical locus, and is
weighted by a \emph{specificity} score in $[0,1]$ recording how exclusively it signals
loyalism: $1.0$ for figures that essentially only mean dynastic loss (\mbox{黍離},
\mbox{採薇}), lower for multivalent ones (\mbox{故國}, \mbox{山河}). The primary analysis uses
the high-specificity subset ($\geq 0.8$) so conclusions do not rest on the multivalent tail;
a sensitivity run uses the full lexicon. The neutral control is a frequency-matched set of
pervasive but non-loyalist canonical allusions (parting willows, the moon, examination
success, wine gatherings), used to test whether any rise is specific to loyalism.

\subsection{Detection and density}
Matching is exact substring matching on traditional-character surface forms, longest form
first so that a four-character figure is not also counted as its two-character substrings,
and each distinct entry is credited once per poem. Text in simplified characters is converted
to traditional before matching. For a poem of $N$ Han characters with high-specificity
loyalist entries $E$ present, the loyalist density is
\[
\rho_{\text{loy}} \;=\; \frac{100}{N}\sum_{e\in E}\, s_e ,
\]
where $s_e$ is the specificity weight; the control density is defined the same way over
control hits. Densities are reported per hundred characters. Where a source mixes prose and
verse (the gazetteers), we first isolate regulated poems with a verse detector that keeps
runs of clauses of a single canonical line length (five or seven characters) and rejects
runs dense in classical-prose function words.

\subsection{Statistics}
The headline contrast is rupture-era density (poems dated $\geq 1894$) against the stable
mid-Qing baseline (1700--1860), the comparison the bimodal shape motivates. Differences are
tested by a bootstrap over poems ($10{,}000$ resamples, seed fixed), reported with a
$95\%$ confidence interval and a two-sided $p$-value; we also report a difference-in-
differences against the control. We give a block bootstrap clustered by source volume as a
conservative check, and we run the contrast under several detector settings to show it does
not depend on one matching choice. All randomised procedures use a fixed seed.

\subsection{Representational probe}
As an independent test of whether the loyalist register is real rather than a property of our
word list, we embed poems with SikuBERT, a BERT model pre-trained on the \mbox{四庫全書} and
therefore innocent of our lexicon. Each poem is labelled loyalist or control by which
register dominates its allusions; we mean-pool the model's final layer and ask whether a
five-fold cross-validated logistic probe separates the two classes above a label-shuffle
null. If the model that never saw the lexicon can recover the distinction, the register is a
property of the language.

\section{Results}
\subsection{A bimodal curve across two ruptures}
Loyalist-allusion density across the dated anthology is not flat and not monotonic. It is
high at the Ming--Zheng founding, collapses through the settled mid-Qing, and rises again
toward the cession (Table~\ref{tab:era}). The two peaks sit at exactly the two moments when
sovereignty over the island changed hands; the trough sits in the long century of Qing rule
between them. A neutral control of canonical non-loyalist allusion shows no such shape.

\begin{table*}[t]
\centering
\footnotesize
\caption{Loyalist vs.\ neutral-control allusion density (per 100 characters) by era in the
dated anthology ($n=997$). The loyalist register peaks at both sovereignty ruptures and
falls quiet in the settled mid-Qing; the control does not.}
\label{tab:era}
\begin{tabular}{lccc}
\toprule
Era (approx.) & poems & loyalist $\rho$ & control $\rho$ \\
\midrule
Ming--Zheng founding ($\leq$1683) & 425 & 0.239 & 0.13 \\
Early/mid Qing (1700--1860)       & 284 & 0.087 & 0.18 \\
Late Qing (1868--1893)            & 111 & 0.255 & 0.23 \\
Cession / rupture ($\geq$1894)    & 177 & 0.268 & 0.18 \\
\bottomrule
\end{tabular}
\end{table*}

The comparison the shape calls for is the rupture era against the stable mid-Qing trough,
not against everything before it (which is diluted by the high founding era). On that
contrast the rupture-era loyalist density is $3.07$ times the mid-Qing density
($0.087 \to 0.268$; bootstrap difference $+0.181$, $95\%$ CI $[0.089, 0.273]$, $p<0.0001$).
The control density is, over the same comparison, essentially unchanged
($0.18 \to 0.18$, $p=0.99$), so the difference-in-differences is $+0.18$: the rise belongs to
loyalist allusion specifically and not to a general increase in allusiveness. The founding
era is elevated to a similar degree ($2.75\times$ the trough), which is the bimodal pattern
stated in numbers.

\subsection{Which figures carry the rupture}
The surge is not spread evenly across the repertoire. Decomposing the rupture-versus-mid-Qing
difference by trope family (Table~\ref{tab:cat}) shows it is carried by the classical canon
of loyalty and refusal-to-serve, with the dynastic-fall topoi second; the local
Koxinga--Dongning vocabulary contributes only marginally, and the explicit cession-diction
hardly at all on its own. In other words, the 1895 poets did not mainly reach for local
references to their own situation; they reached for Bo Yi and Shu Qi, Wen Tianxiang, Tian
Heng, and Su Wu, the inherited language of the loyal subject who will not serve.

\begin{table*}[t]
\centering
\footnotesize
\caption{Per-family decomposition of the rupture-vs-mid-Qing difference (high-specificity
entries). The surge is driven by the loyalty/refusal canon and dynastic-fall topoi, not by
the local Koxinga layer.}
\label{tab:cat}
\begin{tabular}{lccc}
\toprule
Trope family & mid-Qing & rupture & $\Delta$ (95\% CI), $p$ \\
\midrule
Loyalty / refusal (採薇, 文天祥, 田橫, 蘇武) & 0.045 & 0.139 & $+0.094\ [0.025,0.165]$, $p{=}0.007$ \\
Dynastic fall (黍離, 崖山, 銅駝)            & 0.005 & 0.040 & $+0.035\ [0.003,0.069]$, $p{=}0.036$ \\
Ming--Zheng local (延平, 東寧)             & 0.015 & 0.050 & $+0.035\ [-0.003,0.080]$, $p{=}0.089$ \\
Direct diction (割臺, 孤臣)                & 0.000 & 0.017 & $+0.017\ [0.000,0.044]$, $p{=}0.29$ \\
Exile / displacement (屈原, 新亭)          & 0.022 & 0.021 & $-0.001\ [-0.029,0.032]$, $p{=}0.93$ \\
\bottomrule
\end{tabular}
\end{table*}

\subsection{The result does not depend on the detector}
The headline contrast holds, between $2.5$ and $3.1$ times the baseline, under every variant
of the detector we tried: the full lexicon as well as the high-specificity subset, raw token
counts as well as weighted presence, and crucially with the local Koxinga layer removed, with
the explicit cession-diction removed, and with both removed at once
(Table~\ref{tab:robust}). Even stripped of both the locally-specific layers, the rupture era
runs $2.79$ times the mid-Qing trough ($p=0.003$), which is the strongest evidence that the
signal lives in the deep classical repertoire rather than in a handful of topical words.

\begin{table*}[t]
\centering
\footnotesize
\caption{Robustness of the rupture-vs-mid-Qing contrast across detector settings. Significant
and $2.5$--$3.1\times$ throughout.}
\label{tab:robust}
\begin{tabular}{lccc}
\toprule
Detector setting & ratio & $\Delta$ & $p$ \\
\midrule
High-specificity ($\geq 0.8$), presence & $3.07\times$ & $0.181$ & $<0.0001$ \\
Full lexicon, presence                  & $2.74\times$ & $0.258$ & $<0.0001$ \\
Full lexicon, raw tokens                & $2.48\times$ & $0.340$ & $<0.0001$ \\
Minus Ming--Zheng layer                 & $3.02\times$ & $0.146$ & $0.0005$ \\
Minus direct diction                    & $2.88\times$ & $0.164$ & $<0.0001$ \\
Minus both local layers                 & $2.79\times$ & $0.129$ & $0.003$ \\
\bottomrule
\end{tabular}
\end{table*}

\subsection{An independent corpus reproduces the baseline}
The diachronic curve rests on one editor's anthology, which raises the fair worry that the
shape reflects \mbox{連橫}'s selection rather than the literature. We test the worry where it
can be tested. The mid-Qing baseline of the anthology should, if it is real and not an
artefact of selection, agree with an independently compiled corpus of the same period. It
does: the 867 gazetteer verses, assembled by Qing officials, give a mid-Qing loyalist density
of $0.078$, statistically indistinguishable from the anthology's mid-Qing $0.087$ (bootstrap
difference $0.010$, $95\%$ CI $[-0.036, 0.058]$, $p=0.71$). Against this independent baseline
the anthology's rupture poems remain elevated by $+0.191$ ($p<0.0001$) with the control flat
($p=0.32$). The trough, in short, is not something \mbox{連橫} manufactured.

\subsection{The same shape inside one poet's life}
The clearest view of the rupture as a lived event comes from a single uncurated collection.
\mbox{許南英}'s chronologically arranged works move from a pre-rupture baseline of $0.102$ to
$0.131$ in the cession window of 1895--1900 and back down to $0.080$ after 1900, and an
unsupervised changepoint on the series falls at 1897, within the rupture. The spike and the
subsequent decay, the ``half-life'' of the title, appear here in a corpus no anthologist
shaped, though with the small numbers of a single life they are directional rather than
decisive. This single-author trajectory and the population question it raises is exactly what
the unreachable national corpus would let us pursue at scale.

\subsection{A model that never saw the lexicon finds the register}
If the loyalist register were only an artefact of our word list, a language model with no
knowledge of that list should not be able to tell loyalist poems from control poems. It can.
Embedding the anthology poems with SikuBERT and training a cross-validated logistic probe to
separate the two classes gives an accuracy of $0.70$ against a label-shuffle null of $0.51$
($p<0.0001$; Table~\ref{tab:probe}). The separation is reliably decodable rather than
visually clean (the silhouette is small), which is the honest description of an effect that a
classifier recovers but that does not split the space into tidy clusters. Restricting to
rupture-era poems the probe reaches $0.81$, and on mid-Qing poems $0.74$; both clear their
nulls, and the register is, if anything, a little sharper at the rupture.

\begin{table*}[t]
\centering
\footnotesize
\caption{Representational probe with SikuBERT (five-fold CV accuracy vs.\ label-shuffle null).
The register is decodable by a model innocent of our lexicon.}
\label{tab:probe}
\begin{tabular}{lcccc}
\toprule
Subset & poems & CV accuracy & null & $p$ \\
\midrule
All eras           & 200 & 0.700 & 0.511 & $<0.0001$ \\
Rupture ($\geq$1894) & 43  & 0.808 & 0.567 & $<0.0001$ \\
Mid-Qing (1700--1860) & 42 & 0.736 & 0.581 & $0.02$ \\
\bottomrule
\end{tabular}
\end{table*}

\section{What the data say, and what they do not}
Read together, the analyses tell a consistent story. The language of dynastic loss is loud at
both moments when Taiwan's sovereignty changed and quiet in the century between; the rise at
1895 is specific to loyalist allusion and not to allusion in general; it is carried by the
inherited canon of loyal refusal rather than by local references; it survives every change of
detector; an independently compiled corpus reproduces the baseline it is measured against;
and a language model with no access to the lexicon recovers the same distinction. That is a
considerable amount of mutually independent support for a claim that has, until now, rested on
close reading alone.

We are equally clear about the limits. A poet-level check, comparing poets whom historians
label loyalist against the modernist \mbox{賴和}, points the right way but is not significant:
\mbox{丘逢甲}'s own works are markedly loyalist ($0.177$), yet pooled across loyalist poets the
lyric signal is concentrated in him rather than uniform ($0.091$ vs.\ $0.067$, $p\approx0.24$).
More importantly, the diachronic contrast that is decisive at the level of individual poems
becomes non-significant when we cluster conservatively by the six anthology volumes
($p=0.12$); volume clustering is coarse and partly confounded with the time axis itself, since
the volumes are ordered chronologically, but we report it rather than hide it. The conclusions
that do not depend on that single clustering choice, the per-family decomposition, the
independent gazetteer baseline, and the model probe, are the ones on which we rest the
argument.

\section{Limitations}
The decisive remaining question is editorial selection at the rupture peak. We have closed it
for the baseline, by independent replication, and shown the peak is robust against that
independent baseline, but the peak's own significance still draws on one anthology. The corpus
that would settle it, the national \mbox{全臺詩} of some fifty thousand dated poems across many
independent hands, is not currently reachable online and sits otherwise in print and behind a
subscription; our pipeline is built to run on it unchanged. Beyond this, dating by editorial
marker is coarse, at the granularity of reign eras rather than years; lexicon matching favours
precision over recall and so undercounts paraphrased or oblique allusion; and the gazetteer
baseline, while independent, is official descriptive verse and so differs in register from
personal lyric, which is why we use it to corroborate the baseline rather than to test the
peak directly. None of these touches the internal logic of the comparisons; each marks a place
where more or better data would sharpen the estimate.

\section{Conclusion}
We set out to ask whether the second fall of a realm, the cession of Taiwan in 1895, left a
countable trace in the language of the island's classical poetry, and whether that trace
echoed the language of the first fall that had made Taiwan a loyalist refuge two centuries
before. The answer the reachable data give is yes, with the bimodal shape that the hypothesis
predicts: loyalist allusion peaks at both ruptures and falls quiet between them, the rise is
specific to loyalism, it is carried by the deep canon of loyal refusal, and it is recoverable
by a model that never saw our word list. The one analysis that would make the case airtight,
the full population in place of a single anthology, is specified and waiting on access. What we
offer in the meantime is not a final word but a measured one: a falsifiable, openly auditable
account of a literature reaching, at the moment its world changed, for an old language of loss.

\section*{Reproducibility}
All code, the graded loyalist lexicon and neutral control, the dated corpora harvested from
public sources, the verse detector, every statistical procedure (bootstrap intervals,
block bootstrap, permutation tests, the per-family decomposition), and the SikuBERT probe run
deterministically with fixed seeds. The complete repository is available at
\href{https://github.com/aryan35790jp/taiwan-loyalist-allusion}{\texttt{github.com/aryan35790jp/taiwan-loyalist-allusion}}.
The pipeline includes an ingestion path for the national
\mbox{全臺詩} export so that the population analysis can be run without code changes once the
corpus is available.

\begin{thebibliography}{99}
\footnotesize
\bibitem{tsai2014} Chien-hsin Tsai. \emph{A Passage to China: Literature, Loyalism, and
Colonial Taiwan}. Harvard University Asia Center, 2014.
\bibitem{nmtl} National Museum of Taiwan Literature. \emph{全臺詩 (Complete Taiwanese Poems)}.
Tainan: NMTL (multi-volume; digital knowledge base, formerly at xdcm.nmtl.gov.tw).
\bibitem{lianheng} 連橫 (Lian Heng). \emph{臺灣詩乘 (Taiwan Poetry Records)}. c.\ 1921.
\bibitem{xunanying} 許南英 (Xu Nanying). \emph{窺園留草}. Reprinted Taipei, 20th c.
\bibitem{unihan} The Unicode Consortium. \emph{Unicode Han Database (Unihan)}, UAX \#38.
\bibitem{rousseeuw1987} Peter J. Rousseeuw. Silhouettes: a graphical aid to the
interpretation and validation of cluster analysis. \emph{J. Comput. Appl. Math.},
20:53--65, 1987.
\bibitem{efron1979} Bradley Efron. Bootstrap methods: another look at the jackknife.
\emph{The Annals of Statistics}, 7(1):1--26, 1979.
\bibitem{truong2020} Charles Truong, Laurent Oudre, and Nicolas Vayatis. Selective review of
offline change point detection methods. \emph{Signal Processing}, 167:107299, 2020.
\bibitem{sikubert} Dongbo Wang et al. SikuBERT: pre-trained language models for digital
humanities research in classical Chinese. 2021.
\bibitem{opencc} BYVoid and contributors. \emph{OpenCC: Open Chinese Convert}.
\url{https://github.com/BYVoid/OpenCC}.
\bibitem{coffee2018} Neil Coffee et al. Intertextuality in the digital age and the
computational study of allusion. \emph{Digital scholarship in the humanities}, 2018.
\end{thebibliography}

\end{document}
